\chapter{Инструкция по сборке и запуску}
\label{app:build}

В данном приложении описаны шаги для сборки и запуска программы.

\section{Требования}

\subsection{Системные требования}
\begin{itemize}
    \item Операционная система: Linux, macOS или Windows (с WSL)
    \item Компилятор с поддержкой C++20 (GCC 10+, Clang 12+)
    \item CMake версии 3.16 или выше
    \item Python 3.8+ (для визуализатора)
\end{itemize}

\subsection{Зависимости C++}
Проект не имеет внешних зависимостей — используется только стандартная библиотека C++20.

\subsection{Зависимости Python (для визуализатора)}
\begin{verbatim}
pip install plotly pandas numpy
\end{verbatim}

\section{Получение исходного кода}

% TODO: Раскомментировать после публикации репозитория
% \begin{verbatim}
% git clone https://github.com/Straple/Palletizer.git
% cd Palletizer
% \end{verbatim}

\section{Сборка проекта}

\subsection{Базовая сборка}
\begin{verbatim}
mkdir build
cd build
cmake ..
make -j$(nproc)
\end{verbatim}

\subsection{Сборка с указанием числа потоков}
\begin{verbatim}
cmake -DTHREADS_NUM=32 ..
make -j$(nproc)
\end{verbatim}

\subsection{Сборка в режиме Release}
\begin{verbatim}
cmake -DCMAKE_BUILD_TYPE=Release ..
make -j$(nproc)
\end{verbatim}

\section{Запуск программы}

\subsection{Основная программа}
\begin{verbatim}
./palletizer <входной_файл.csv> [выходной_файл.csv]
\end{verbatim}

\subsection{Пример запуска}
\begin{verbatim}
./palletizer ../tests/100.csv answer.csv
\end{verbatim}

\section{Запуск визуализатора}

\begin{verbatim}
python3 src/scripts/visualizer.py answer.csv \
    --pallet-length 1200 \
    --pallet-width 800 \
    --color-by sku \
    --opacity 0.4 \
    --out visualization.html
\end{verbatim}

После выполнения команды откроется браузер с интерактивной 3D-визуализацией.

\section{Структура проекта}

\begin{verbatim}
Palletizer/
├── CMakeLists.txt          # Конфигурация сборки
├── src/
│   ├── main.cpp            # Точка входа
│   ├── settings.hpp        # Настройки
│   ├── objects/            # Структуры данных
│   │   ├── test_data.*     # Входные данные
│   │   ├── answer.*        # Результат укладки
│   │   └── metrics.*       # Метрики
│   ├── solvers/            # Алгоритмы
│   │   ├── height_handler* # Карта высот (5 реализаций)
│   │   ├── greedy/         # GreedySolver
│   │   └── lns/            # LNSSolver
│   ├── utils/              # Утилиты
│   └── scripts/            # Python-скрипты
│       └── visualizer.py   # 3D-визуализатор
├── tests/                  # Тестовые данные (436 файлов)
└── papers/                 # Документация (LaTeX)
\end{verbatim}

% TODO: Добавить диаграмму классов (figures/class_diagram.png)
% Показать: основные классы (Box, TestData, Answer, HeightHandler, Solver)
% и связи между ними

\section{Запуск тестов}

\subsection{Тестирование HeightHandler}
\begin{verbatim}
./test_height_handler
\end{verbatim}

Выводит сравнительную таблицу производительности всех реализаций HeightHandler.

\subsection{Обучение весов мутаций LNS}
\begin{verbatim}
./train_lns_solver_mutations
\end{verbatim}

Запускает процесс оптимизации весов операторов мутации.

\section{Формат входных данных}

CSV-файл с заголовком:
\begin{verbatim}
SKU,Quantity,Length,Width,Height,Weight,Strength,Aisle,Caustic
900001,6,600,400,200,12000,4,1,0
900002,5,200,400,200,5000,4,2,0
\end{verbatim}

\section{Формат выходных данных}

CSV-файл с координатами уложенных коробок:
\begin{verbatim}
SKU,x,y,z,X,Y,Z,Aisle,Weight
900001,0,0,0,600,400,200,1,12000
900001,600,0,0,1200,400,200,1,12000
\end{verbatim}

где $(x, y, z)$ --- нижний угол коробки, $(X, Y, Z)$ --- верхний угол.
